%引入纯矢量角标定义=============================================================
\input{tikz_assets/box_corner_mark.tex}

%通用盒子定义===============================================================================
\newlength{\textwd} %自动计算文本长度

% 通用三栏标题盒命令1（左右图片 + 中间自适应文字）
\newcommand{\TitleBoxI}[5]{%
	\setlength{\textwd}{0pt}%
	\sbox0{\fontsize{20}{22}\selectfont\SiYuan\XeTeXinterchartokenstate=0 #2}%
	\setlength{\textwd}{\widthof{\usebox0}}%
	\begin{tabular}[c]{@{}c@{}c@{}c@{}}
		\raisebox{-\height + 12pt}{%
			\resizebox{!}{24pt}{\input{#1}}%
		} &
		\raisebox{-\height + 12pt}{%
			\colorbox{#4}{%
				\parbox[c][18pt][c]{\textwd}{%
					\centering
					\fontsize{18}{22}\selectfont\SiYuan\XeTeXinterchartokenstate=0\color{#5} #2%
				}%
			}%
		} &
		\raisebox{-\height + 12pt}{%
			\resizebox{!}{24pt}{\input{#3}}%
		}
	\end{tabular}%
}

% 通用三栏标题盒命令2（左右图片 + 中间自适应文字）
\newcommand{\TitleBoxII}[5]{%
	\setlength{\textwd}{0pt}%
	\sbox0{\fontsize{16}{22}\selectfont\SiYuan\XeTeXinterchartokenstate=0 #2}%
	\setlength{\textwd}{\widthof{\usebox0}}%
	\begin{tabular}[c]{@{}c@{}c@{}c@{}}
		\raisebox{-\height + 12pt}{%
			\resizebox{!}{24pt}{\input{#1}}%
		} &
		\raisebox{-\height + 12pt}{%
			\colorbox{#4}{%
				\parbox[c][18pt][c]{\textwd}{%
					\centering
					\fontsize{14}{22}\selectfont\SiYuan\XeTeXinterchartokenstate=0\color{#5} #2%
					\vspace{-4pt}
				}%
			}%
		} &
		\raisebox{-\height + 12pt}{%
			\resizebox{!}{24pt}{\input{#3}}%
		}
	\end{tabular}%
}

% 通用三栏标题盒命令3（左右图片 + 中间自适应文字）
\newcommand{\TitleBoxIII}[5]{%
	\setlength{\textwd}{0pt}%
	\sbox0{\fontsize{22}{22}\selectfont\SiYuan\XeTeXinterchartokenstate=0 #2}%
	\setlength{\textwd}{\widthof{\usebox0}}%
	\begin{tabular}[c]{@{}c@{}c@{}c@{}}
		\raisebox{-\height + 12pt}{%
			\resizebox{!}{24pt}{\input{#1}}%
		} &
		\raisebox{-\height + 12pt}{%
			\colorbox{#4}{%
				\parbox[c][18pt][c]{\textwd}{%
					\centering
					\fontsize{14}{22}\selectfont\SiYuan\XeTeXinterchartokenstate=0\color{#5} #2%
					\vspace{0pt}
				}%
			}%
		} &
		\raisebox{-\height + 12pt}{%
			\resizebox{!}{24pt}{\input{#3}}%
		}
	\end{tabular}%
}

% 通用三栏标题盒命令4（左右图片 + 中间自适应文字）
\newcommand{\TitleBoxIV}[5]{%
	\setlength{\textwd}{0pt}%
	\sbox0{\fontsize{16}{22}\selectfont\SiYuan\XeTeXinterchartokenstate=0 #2}%
	\setlength{\textwd}{\widthof{\usebox0}}%
	\begin{tabular}[c]{@{}c@{}c@{}c@{}}
		\raisebox{-\height + 12pt}{%
			\resizebox{!}{24pt}{\input{#1}}%
		} &
		\raisebox{-\height + 12pt}{%
			\colorbox{#4}{%
				\parbox[c][18pt][c]{\textwd}{%
					\centering
					\fontsize{14}{22}\selectfont\SiYuan\XeTeXinterchartokenstate=0\color{#5} #2%
					\vspace{0pt}
				}%
			}%
		} &
		\raisebox{-\height + 12pt}{%
			\resizebox{!}{24pt}{\input{#3}}%
		}
	\end{tabular}%
}

%相关盒子定义============================================================
% 1. 一级标题和二级标题盒子设置
\titleformat{\subsection}
{\raggedright}{}{0pt}{
	\TitleBoxI{tikz_assets/box1_1.tex}{#1}{tikz_assets/box1_2.tex}{mibai}{qiancheng}
	\vspace{8pt}
}

% 2. 相关数学内容盒子设置
\newtcolorbox{definition}[1]{
	enhanced,
	sharp corners,
	colback=qianbai,
	colframe=definition,
	boxrule=2.5pt,
	fontupper={\kaishu},
	halign=justify,
	breakable,
	title={\TitleBoxII{tikz_assets/definition_1.tex}{#1}{tikz_assets/definition_2.tex}{heise}{mibai}},
	attach boxed title to top left={yshift=-2.5pt, xshift=0mm},
	boxed title style={
		boxrule=0pt,
		boxsep=0pt,
		top=-2.1pt,
		bottom=0pt,
		left=0pt,
		right=0pt,
		colback=qianbai,
		baseline=\ht\strutbox
	},
	after upper={\hfill\resizebox{!}{8pt}{\BoxCornerMark[definition]}},
}

\newtcolorbox{theorem}[1]{
	enhanced,
	sharp corners,
	colback=qianbai,
	colframe=theorem,
	boxrule=2.5pt,
	fontupper={\kaishu},
	halign=justify,
	breakable,
	title={\TitleBoxII{tikz_assets/theorem_1.tex}{#1}{tikz_assets/theorem_2.tex}{heise}{mibai}},
	attach boxed title to top left={yshift=-2.5pt, xshift=0mm},
	boxed title style={
		boxrule=0pt,
		boxsep=0pt,
		top=-2.1pt,
		bottom=0pt,
		left=0pt,
		right=0pt,
		colback=qianbai,
		baseline=\ht\strutbox
	},
	after upper={\hfill\resizebox{!}{8pt}{\BoxCornerMark[theorem]}},
}

\newtcolorbox{proposition}[1]{
	enhanced,
	sharp corners,
	colback=qianbai,
	colframe=proposition,
	boxrule=2.5pt,
	fontupper={\kaishu},
	halign=justify,
	breakable,
	title={\TitleBoxII{tikz_assets/proposition_1.tex}{#1}{tikz_assets/proposition_2.tex}{heise}{mibai}},
	attach boxed title to top left={yshift=-2.5pt, xshift=0mm},
	boxed title style={
		boxrule=0pt,
		boxsep=0pt,
		top=-2.1pt,
		bottom=0pt,
		left=0pt,
		right=0pt,
		colback=qianbai,
		baseline=\ht\strutbox
	},
	after upper={\hfill\resizebox{!}{8pt}{\BoxCornerMark[proposition]}},
}

\newtcolorbox{corollary}[1]{
	enhanced,
	sharp corners,
	colback=qianbai,
	colframe=corollary,
	boxrule=2.5pt,
	fontupper={\kaishu},
	halign=justify,
	breakable,
	title={\TitleBoxII{tikz_assets/corollary_1.tex}{#1}{tikz_assets/corollary_2.tex}{heise}{mibai}},
	attach boxed title to top left={yshift=-2.5pt, xshift=0mm},
	boxed title style={
		boxrule=0pt,
		boxsep=0pt,
		top=-2.1pt,
		bottom=0pt,
		left=0pt,
		right=0pt,
		colback=qianbai,
		baseline=\ht\strutbox
	},
	after upper={\hfill\resizebox{!}{8pt}{\BoxCornerMark[corollary]}},
}

\newtcolorbox{lemma}[1]{
	enhanced,
	sharp corners,
	colback=qianbai,
	colframe=lemma,
	boxrule=2.5pt,
	fontupper={\kaishu},
	halign=justify,
	breakable,
	title={\TitleBoxII{tikz_assets/lemma_1.tex}{#1}{tikz_assets/lemma_2.tex}{heise}{mibai}},
	attach boxed title to top left={yshift=-2.5pt, xshift=0mm},
	boxed title style={
		boxrule=0pt,
		boxsep=0pt,
		top=-2.1pt,
		bottom=0pt,
		left=0pt,
		right=0pt,
		colback=qianbai,
		baseline=\ht\strutbox
	},
	after upper={\hfill\resizebox{!}{8pt}{\BoxCornerMark[lemma]}},
}

\newtcolorbox{example}[1]{
	enhanced,
	sharp corners,
	colback=qianbai,
	colframe=example,
	boxrule=2.5pt,
	fontupper={\kaishu},
	halign=justify,
	breakable,
	title={\TitleBoxII{tikz_assets/example_1.tex}{#1}{tikz_assets/example_2.tex}{heise}{mibai}},
	attach boxed title to top left={yshift=-2.5pt, xshift=0mm},
	boxed title style={
		boxrule=0pt,
		boxsep=0pt,
		top=-2.1pt,
		bottom=0pt,
		left=0pt,
		right=0pt,
		colback=qianbai,
		baseline=\ht\strutbox
	},
	after upper={\hfill\resizebox{!}{8pt}{\BoxCornerMark[example]}},
}

\newtcolorbox{exercise}[1]{
	enhanced,
	sharp corners,
	colback=qianbai,
	colframe=qiancheng,
	boxrule=2.5pt,
	fontupper={\kaishu},
	halign=justify,
	breakable,
	title={\TitleBoxII{tikz_assets/exercise_1.tex}{#1}{tikz_assets/exercise_2.tex}{mibai}{qiancheng}},
	attach boxed title to top left={yshift=0pt, xshift=0mm},
	boxed title style={
		boxrule=0pt,
		boxsep=0pt,
		top=-2.1pt,
		bottom=0pt,
		left=0pt,
		right=0pt,
		colback=qianbai,
		baseline=\ht\strutbox
	},
}

\newtcolorbox{question}[1]{
	enhanced,
	sharp corners,
	colback=qianbai,
	colframe=qiancheng,
	boxrule=2.5pt,
	fontupper={\kaishu},
	halign=justify,
	breakable,
	title={\TitleBoxII{tikz_assets/exercise_1.tex}{#1}{tikz_assets/question_2.tex}{mibai}{qiancheng}},
	attach boxed title to top left={yshift=0pt, xshift=0mm},
	boxed title style={
		boxrule=0pt,
		boxsep=0pt,
		top=-2.1pt,
		bottom=0pt,
		left=0pt,
		right=0pt,
		colback=qianbai,
		baseline=\ht\strutbox
	},
}

\newtcolorbox{prove}[1][]{
	enhanced,
	sharp corners,
	colback=qianbai,  
	colframe=qiancheng,  
	fontupper={\kaishu}, 
	halign=justify,
	leftrule=5pt, 
	rightrule=0pt,      
	toprule=0pt,        
	bottomrule=0pt,     
	breakable,
	before upper={
		\setstretch{1.5}\noindent
		\textcolor{qiancheng}{\XeTeXinterchartokenstate=0\SiYuan\textbf{证明：}} 
	},
	after upper={\hfill\textcolor{qiancheng}{$\blacksquare$}},
	#1 
}

\newtcolorbox{remark}[1][]{
	enhanced,
	sharp corners,
	colback=qianbai,  
	colframe=qiancheng,  
	fontupper={\kaishu},  
	halign=justify,
	leftrule=5pt, 
	rightrule=0pt,      
	toprule=0pt,        
	bottomrule=0pt,     
	breakable,
	before upper={
		\setstretch{1.5}\noindent
		\textcolor{qiancheng}{\XeTeXinterchartokenstate=0\SiYuan\textbf{点评：}} 
	},
	after upper={\hfill\textcolor{qiancheng}{$\blacksquare$}},
	#1 
}

% 3. 其他盒子设置
\newtcolorbox{box1}[1]{
	enhanced,
	sharp corners,
	colback=qianbai,
	colframe=qiancheng,
	fontupper={\kaishu},
	halign=justify,
	boxrule=2pt,
	breakable,
	title={\TitleBoxIII{tikz_assets/box1_1.tex}{#1}{tikz_assets/box1_2.tex}{mibai}{qiancheng}},
	attach boxed title to top left={yshift=0pt, xshift=0mm},
	boxed title style={
		boxrule=0pt,
		boxsep=0pt,
		top=-2.1pt,
		bottom=0pt,
		left=0pt,
		right=0pt,
		colback=qianbai,
		baseline=\ht\strutbox
	},
}

\newtcolorbox{box2}[1]{
	enhanced,
	sharp corners,
	colback=qianbai,
	colframe=qiancheng,
	fontupper={\kaishu},
	halign=justify,
	boxrule=2pt,
	breakable,
	title={\TitleBoxIV{tikz_assets/box2_1.tex}{#1}{tikz_assets/box2_2.tex}{qiancheng}{mibai}},
	attach boxed title to top left={yshift=0pt, xshift=0mm},
	boxed title style={
		boxrule=0pt,
		boxsep=0pt,
		top=-2.1pt,
		bottom=0pt,
		left=0pt,
		right=0pt,
		colback=qianbai,
		baseline=\ht\strutbox
	},
}

\newtcolorbox{box3}[1]{
	enhanced,
	sharp corners,
	colback=qianbai,
	colframe=qiancheng,
	fontupper={\kaishu},
	halign=justify,
	boxrule=2pt,
	breakable,
	title={\TitleBoxIII{tikz_assets/box3_1.tex}{#1}{tikz_assets/box3_2.tex}{heise}{mibai}},
	attach boxed title to top left={yshift=0pt, xshift=0mm},
	boxed title style={
		boxrule=0pt,
		boxsep=0pt,
		top=-2.1pt,
		bottom=0pt,
		left=0pt,
		right=0pt,
		colback=qianbai,
		baseline=\ht\strutbox
	},
}

\newtcolorbox{box4}[1]{
	enhanced,
	sharp corners,
	colback=qianbai,
	colframe=qiancheng,
	fontupper={\kaishu},
	halign=justify,
	boxrule=2pt,
	breakable,
	title={\TitleBoxIV{tikz_assets/box4_1.tex}{#1}{tikz_assets/box4_2.tex}{mibai}{qiancheng}},
	attach boxed title to top left={yshift=0pt, xshift=0mm},
	boxed title style={
		boxrule=0pt,
		boxsep=0pt,
		top=-2.1pt,
		bottom=0pt,
		left=0pt,
		right=0pt,
		colback=qianbai,
		baseline=\ht\strutbox
	},
}


%自定义其他命令===================================================
% 1. 自定义分割线
\newcommand{\Ark}[0]{
	\par\vspace{-10pt}
	\noindent
	\resizebox{\linewidth}{2pt}{\input{tikz_assets/separator_1.tex}}
	\par\vspace{-4pt}
}

% 2. 自定义文字居中
\newcommand{\C}[1]{\makebox[\linewidth][c]{#1}}

% 3. 自定义增加间隔
\newcommand{\Q}[0]{\vspace{4pt}}

% 4. 自定义添加文本背景
\newcommand{\B}[1]{%
	\tikz[baseline=(char.base)]{%
		\node[inner sep=0pt, outer sep=0pt] (char) {\phantom{#1}};
		\fill[qiancheng] (char.south west) rectangle ([xshift=0pt, yshift=1.5pt]char.south east);
		\node[inner sep=0pt, outer sep=0pt] (char) {#1};
	}%
}